\documentclass{article}
\usepackage{amsmath, amssymb, geometry}
\geometry{a4paper, margin=1in}

\title{Digest: Lorentz/Biquaternion Equivalence}
\author{InfoGeometry Agent}
\date{\today}

\begin{document}
\maketitle

\section{Introduction}
This digest formalizes the Lorentz/Biquaternion Equivalence, establishing the precise mapping between the restricted Lorentz group $SO^+(1,3)$ and the multiplicative group of unit biquaternions (complexified quaternions $\mathbb{C} \otimes \mathbb{H}$).

\section{Biquaternion Algebraic Structure}
The algebra of biquaternions provides a canonical representation for spacetime transformations. A biquaternion $Q \in \mathbb{C} \otimes \mathbb{H}$ can be written as $Q = q_1 + i q_2$ where $q_1, q_2$ are real quaternions.
The norm of $Q$ is defined via the biquaternion conjugate $Q \bar{Q}$. The group of unit biquaternions ($Q \bar{Q} = 1$) is isomorphic to the spin group $SL(2, \mathbb{C})$.

\section{Lorentz Transformations}
A four-vector $x = (t, x, y, z)$ is represented as a Hermitian biquaternion $X = t + x \sigma_1 + y \sigma_2 + z \sigma_3$.
The action of the Lorentz group is generated by the double-sided transformation:
\begin{equation}
X' = Q X Q^\dagger
\end{equation}
where $Q^\dagger$ is the complex conjugate of the biquaternion conjugate.

\section{Equivalence Theorem}
The Lorentz/Biquaternion Equivalence theorem mathematically asserts that the continuous Lorentz boost generators and spatial rotation generators are exactly spanned by the respective real and imaginary biquaternion basis components, ensuring that spacetime geometry natively reduces to biquaternion algebra.

\end{document}
